\chapter{נוף האיומים הלאומי והחזון האסטרטגי}
\label{ch:vision}

\section{מבוא: עידן התודעה הסינתטית והאתגר הדמוקרטי}
\label{sec:vision_intro}

המרחב הציבורי הדיגיטלי בעשור השלישי של המאה ה-21 עובר תמורה מבנית חסרת תקדים. אם בעשור הקודם התמקד האיום על מערכות דמוקרטיות בהפצה מתואמת של מידע כוזב (\textenglish{Disinformation}) באמצעות רשתות בוטים בסיסיות וחוות טרולים אנושיות, הרי שמהפכת הבינה המלאכותית היוצרת (\textenglish{Generative AI}) שינתה את כללי המשחק מן היסוד~\cite{goldstein2023generative}. כיום, שחקנים מדינתיים, ארגוני פרוקסי וגורמים עוינים מסוגלים לייצר, להתאים אישית ולהפיץ תוכן רב-מודאלי --- טקסט, אודיו, תמונה ווידאו --- ברמת אמינות גבוהה, בעלות שולית הקרובה לאפס ובקצב החורג מיכולת הניטור והתגובה של מוסדות המדינה המסורתיים.

עבור מדינת ישראל, דמוקרטיה פרלמנטרית הפועלת בסביבה ביטחונית מורכבת ומתמשכת, האתגר הוא כפול ומכופל. מצד אחד, המרחב המקוון הישראלי מהווה זירה מרכזית למימוש חופש הביטוי, להשתתפות אזרחית ולשיח פוליטי תוסס; מצד שני, אותו מרחב עצמו חשוף למבצעי השפעה זרים (\textenglish{Foreign Information Manipulation and Interference -- FIMI}), ללוחמה תודעתית בעתות חירום ובמערכות בחירות, ולניסיונות ממוקדים להעמיק שסעים חברתיים ולשחוק את אמון הציבור במוסדות הדמוקרטיים~\cite{bennett2018disinformation}.

מערכת \textenglish{CivicShield AI} נולדה מתוך ההכרה כי הגנה על חוסנה הדמוקרטי של המדינה מחייבת תשתית לאומית מתקדמת, שקופה ומפוקחת, המשלבת יכולות זיהוי מוקדם מבוססות בינה מלאכותית עם מחויבות בלתי מתפשרת לעקרונות שלטון החוק, חופש הביטוי והגנת הפרטיות. דוח מקדים זה מציג את התשתית התאורטית, המשפטית, הארכיטקטונית והתפעולית של המיזם, כבסיס להיערכות למכרז לאומי רב-שנתי.

\section{מיפוי נוף האיומים הלאומי במרחב המידע}
\label{sec:threat_landscape}

ניתוח שיטתי של איומי ההשפעה במרחב המידע הישראלי והבינלאומי מצביע על חמישה וקטורי איום מרכזיים הפועלים לעיתים קרובות במשולב במסגרת מערכות היברידיות~\cite{rid2020active}:

\begin{enumerate}
    \item \textbf{מבצעי השפעה זרים מתואמים (\textenglish{FIMI}):} הפעלת רשתות של חשבונות מזויפים, אתרי חדשות מתחזים (\textenglish{Doppelgänger Sites}) ודמויות סינתטיות במטרה להחדיר נרטיבים עוינים לשיח הציבורי המקומי, לייצר תחושת כאוס, או להשפיע על תהליכי קבלת החלטות ובחירות.
    \item \textbf{תוכן סינתטי וזיופים עמוקים (\textenglish{Deepfakes}):} יצירת הקלטות קוליות, סרטוני וידאו ותמונות מבוימות של נבחרי ציבור, קציני ביטחון או אזרחים, במטרה לייצר משברי אמון מיידיים (``אפקט ההפתעה של אוקטובר'') או לשבש תפקוד של שירותי חירום~\cite{chesney2019deep}.
    \item \textbf{מיקרו-טרגוטינג קוגניטיבי והקצנה רגשית:} שימוש במודלי שפה גדולים (\textenglish{LLMs}) לניתוח קהלי יעד והתאמת מסרים מפלגים בשפה העברית והערבית, תוך ניצול הטיות קוגניטיביות ורגשות של פחד, כעס וניכור~\cite{tucker2018social}.
    \item \textbf{לוחמה תודעתית בזמן אמת במצבי חירום:} הצפת רשתות חברתיות ואפליקציות מסרים בשמועות כוזבות על אירועים ביטחוניים, תקלות בתשתיות קריטיות או הנחיות סותרות של פיקוד העורף, במטרה לייצר בהלה המונית ולפגוע בתפקוד העורף האזרחי.
    \item \textbf{שחיקת אמון מערכתית (``דיבידנד השקרן'' -- \textenglish{Liar's Dividend}):} מעבר להטעיה נקודתית, עצם הידיעה כי כל תוכן עלול להיות מזויף מאפשרת לשחקנים עוינים להכחיש ראיות אמיתיות ולערער את האמון הבסיסי בעובדות, בתקשורת החופשית ובמוסדות המדינה~\cite{chesney2019deep}.
\end{enumerate}

\subsection{טקסונומיה השוואתית של איומי מידע}

טבלה~\ref{tab:threat_taxonomy} מסכמת את מאפייני האיומים המרכזיים, את רמת הסיכון המוערכת למערכת הדמוקרטית בישראל, ואת חלון הזמן הנדרש לזיהוי ותגובה אפקטיבית.

\begin{table}[htbp]
\centering
\caption{טקסונומיה של איומי מידע לאומיים ודרישות תגובה}
\label{tab:threat_taxonomy}
\begin{tabularx}{\textwidth}{p{3.2cm} X p{2.4cm} p{2.4cm}}
\toprule
\textbf{קטגוריית איום} & \textbf{מנגנון פעולה מרכזי} & \textbf{חלון תגובה נדרש} & \textbf{רמת סיכון לאומית} \\
\midrule
רשתות השפעה זרות (\textenglish{CIB}) & תיאום לא-אותנטי של מאות עד אלפי חשבונות להגברת נרטיבים מפלגים & שעות עד ימים & גבוהה מאוד \\
\addlinespace
זיופים עמוקים (\textenglish{Deepfakes}) & מדיה סינתטית המתחזה למקבלי החלטות או לאירועי אמת & דקות עד שעות & קריטית (בבחירות/חירום) \\
\addlinespace
שיבוש מידע בחירום & הפצת הנחיות כוזבות ושמועות בעת מתקפה או אסון & דקות ספורות & קריטית \\
\addlinespace
אתרי חדשות מתחזים & שכפול עיצוב של כלי תקשורת מוכרים להלבנת מידע כוזב & שעות & גבוהה \\
\addlinespace
הקצנה אלגוריתמית & ניצול מנגנוני המלצה של פלטפורמות להדהוד שיח שנאה & ימים עד שבועות & בינונית--גבוהה \\
\bottomrule
\end{tabularx}
\end{table}

\section{כשלי המצב הקיים ופער היכולות הלאומי}
\label{sec:capability_gap}

בחינת אופן ההתמודדות הנוכחי של גופי המדינה עם איומי מידע חושפת ארבעה פערים מבניים מרכזיים:

\begin{itemize}
    \item \textbf{פיצול ארגוני ומידע לא משותף:} האחריות לניטור איומי מידע מבוזרת בין גופים שונים (מערך הסייבר הלאומי, ועדת הבחירות המרכזית, משרד המשפטים, גופי ביטחון ומערך ההסברה הלאומי), ללא תמונת מצב משותפת (\textenglish{Common Operating Picture}) בזמן אמת.
    \item \textbf{נחיתות טכנולוגית בשפה העברית והערבית:} רוב כלי הניטור המסחריים הבינלאומיים מאומנים בעיקר על אנגלית ואינם מזהים היטב ניואנסים תרבותיים, סלנג מקומי, קודים פוליטיים וערבוב שפות (\textenglish{Code-Switching}) האופייניים לשיח הישראלי.
    \item \textbf{תגובתיות מאוחרת:} בדיקת עובדות אנושית (\textenglish{Fact-Checking}) חיונית אך איטית מטבעה; עד לפרסום הפרכה רשמית, הנרטיב הכוזב כבר צבר תפוצה רחבה והתקבע בתודעה הציבורית~\cite{vosoughi2018spread}.
    \item \textbf{חוסר בהירות נורמטיבית ואמון ציבורי:} היעדר מסגרת שקופה ומבוקרת לפעילות המדינה במרחב המידע מעורר חשש מוצדק מפני ``משרד אמת'' או פגיעה בשיח פוליטי לגיטימי, ובכך פוגע בלגיטימציה הציבורית של מאמצי ההגנה.
\end{itemize}

\section{החזון האסטרטגי של \textenglish{CivicShield AI}}
\label{sec:strategic_vision}

חזון \textenglish{CivicShield AI} הוא הקמת תשתית לאומית חכמה, מבוזרת ושקופה להגנה על יושרת המידע (\textenglish{Information Integrity}) ועל החוסן הדמוקרטי בישראל. המערכת אינה מבקשת ``לצנזר'' דעות או להכריע בוויכוחים פוליטיים לגיטימיים; תפקידה הוא לזהות מניפולציות התנהגותיות לא-אותנטיות, תוכן סינתטי מטעה ומבצעי השפעה זרים, ולספק למקבלי ההחלטות, לגופי הבחירות ולציבור הרחב כלים אמינים, מבוססי ראיות ומהימנים לקבלת החלטות מושכלת.

\subsection{עקרונות הליבה האסטרטגיים}

האסטרטגיה של \textenglish{CivicShield AI} נשענת על חמישה עקרונות יסוד בלתי ניתנים לפשרה:

\begin{enumerate}
    \item \textbf{התמקדות בהתנהגות ולא בדעה (\textenglish{Behavior, Not Speech}):} המערכת מנתחת דפוסי תיאום לא-אותנטיים (\textenglish{Coordinated Inauthentic Behavior}), מקוריות מדיה (\textenglish{Provenance}) וטביעות אצבע טכנולוגיות של אוטומציה --- ולא את העמדה הפוליטית או האידאולוגית של הדובר.
    \item \textbf{בינה מלאכותית בת-הסבר עם אדם בלולאה (\textenglish{Human-in-the-Loop XAI}):} שום פעולה אופרטיבית, התרעה רשמית או פנייה לפלטפורמה אינה מתבצעת באופן אוטונומי מלא. כל ממצא אלגוריתמי מלווה בהסבר ברור, בציון ודאות ובראיות ניתנות לביקורת אנושית.
    \item \textbf{פרטיות מובנית (\textenglish{Privacy by Design}):} המערכת פועלת אך ורק על מידע פומבי (\textenglish{OSINT}) תוך אנונימיזציה אגרגטיבית של אזרחים פרטיים, ללא מעקב אישי וללא אגירת פרופילים פוליטיים.
    \item \textbf{ריבונות טכנולוגית ולשונית:} פיתוח מודלים ייעודיים לעברית, ערבית, רוסית ואנגלית, המאוחסנים בתשתית ענן ממשלתית מאובטחת (``נימבוס'') תחת שליטה לאומית מלאה.
    \item \textbf{שקיפות ופיקוח דמוקרטי רב-שכבתי:} פעילות המערכת כפופה לביקורת חיצונית בלתי תלויה של ועדה ציבורית, מבקר המדינה וגופי הגנת הפרטיות, כולל פרסום דוחות שקיפות תקופתיים לציבור.
\end{enumerate}

\section{מפת דרכים תלת-שלבית ויעדים מדידים}
\label{sec:roadmap}

הקמת \textenglish{CivicShield AI} מתוכננת כמערכת מדורגת הפרוסה על פני שלוש שנים (36 חודשים), המאפשרת הפקת ערך מבצעי מוקדם לצד בניית אמון ציבורי ובקרה משפטית הדוקה:

\begin{itemize}
    \item \textbf{שלב א' (חודשים 1--12): תשתית ליבה וזיהוי מוקדם.} הקמת מנוע ניטור פומבי רב-לשוני, פיתוח גלאי זיופים עמוקים (\textenglish{Deepfake Detection}) ותשתית זיהוי רשתות מתואמות, והפעלת פיילוט מבוקר מול ועדת הבחירות המרכזית ומערך הסייבר.
    \item \textbf{שלב ב' (חודשים 13--24): שילוב מערכתי וחיסון קוגניטיבי.} חיבור לפורטל התרעות לאומי בזמן אמת, שיתופי פעולה עם כלי תקשורת וארגוני בדיקת עובדות, והטמעת מנגנוני ``הקדמת תרופה למכה'' (\textenglish{Prebunking}) בשיתוף מערכת החינוך והחברה האזרחית.
    \item \textbf{שלב ג' (חודשים 25--36): חוסן לאומי מלא ושיתוף פעולה בינלאומי.} הרחבה לתקני אימות מקוריות מדיה (\textenglish{C2PA}), ממשקי שיתוף מודיעין השפעה עם דמוקרטיות שותפות, וביסוס מנגנון ביקורת אלגוריתמית מתמיד (\textenglish{Continuous Red-Teaming}).
\end{itemize}

\subsection{מדדי הצלחה מרכזיים (\textenglish{KPIs})}

טבלה~\ref{tab:strategic_kpis} מציגה את יעדי הביצוע הכמותיים של המערכת לאורך שלבי ההטמעה, המהווים בסיס למנגנון התשלום מבוסס-התוצאות במכרז הממשלתי.

\begin{table}[htbp]
\centering
\caption{מדדי הצלחה מרכזיים (\textenglish{KPIs}) למיזם \textenglish{CivicShield AI}}
\label{tab:strategic_kpis}
\begin{tabularx}{\textwidth}{p{3.8cm} X p{2.2cm} p{2.2cm}}
\toprule
\textbf{מדד ביצוע} & \textbf{הגדרה תפעולית} & \textbf{יעד שנה 1} & \textbf{יעד שנה 3} \\
\midrule
זמן ממוצע לזיהוי (\textenglish{MTTD}) & הזמן מתחילת קמפיין השפעה מתואם ועד להפקת התרעה מאומתת & $< 4$ שעות & $< 30$ דקות \\
\addlinespace
דיוק זיהוי (\textenglish{Precision / F1}) & שיעור דיוק בזיהוי רשתות לא-אותנטיות ומדיה סינתטית בעברית/ערבית & $F_1 \ge 0.85$ & $F_1 \ge 0.93$ \\
\addlinespace
שיעור התרעות שווא (\textenglish{FPR}) & אחוז תוכן אזרחי לגיטימי שסומן בטעות כמניפולציה מתואמת & $< 2.0\%$ & $< 0.5\%$ \\
\addlinespace
כיסוי שקיפות וביקורת & אחוז ההחלטות האופרטיביות המלוות בתיעוד מלא וביקורת אנושית & $100\%$ & $100\%$ \\
\bottomrule
\end{tabularx}
\end{table}

\section{סיכום פרק}
\label{sec:vision_summary}

הגנה על המרחב הדמוקרטי בעידן הבינה המלאכותית אינה משימה טכנולוגית בלבד, אלא אתגר לאומי רב-תחומי המחייב סינתזה בין הנדסה מתקדמת, מדעי החברה והמדינה ומשפט חוקתי. הפרק הבא מניח את התשתית התאורטית והמשפטית הנדרשת כדי להבטיח כי \textenglish{CivicShield AI} תפעל לא רק ביעילות מבצעית גבוהה, אלא גם כעוגן של לגיטימציה דמוקרטית וזכויות אדם.
